\section{Apéndice: solución al ABC 481 F}

\begin{frame}
  \vspace{1.5cm}
  \centering
  {\LARGE Apéndice: solución al ABC 481 F}
\end{frame}

\begin{frame}{El orden de los costos}
  Un camino con más aristas tiene más dígitos.
  Con dígitos del 1 al 9, ese número es mayor.
  Repetir un nodo agrega dígitos, así que no mejora el mínimo.

  \medskip
  Entre dos caminos con la misma cantidad de aristas, el menor número es el menor en orden lexicográfico.

  \medskip
  Alcanza con buscar el camino lexicográficamente menor entre los de longitud mínima.
\end{frame}

\begin{frame}{Capa por capa}
  El nodo 1 empieza con la string vacía, id \(0\).
  En la capa \(k\) se extienden los caminos de la capa \(k-1\) con una arista, y cada nodo se queda con el menor par \((\mathrm{id}, d)\).
  Esos pares ganadores se comprimen, y los ids nuevos se usan en la capa siguiente.

  \vspace{0.45em}
  \begin{columns}[c,onlytextwidth]
    \begin{column}{0.34\textwidth}
      \centering
      \begin{tikzpicture}[scale=0.86]
        \node[v] (a) at (0,0) {\(1\)};
        \node[v] (b) at (-1.35,-1.45) {\(2\)};
        \node[v] (c) at (1.35,-1.45) {\(3\)};
        \node[v] (d) at (0,-2.9) {\(4\)};
        \draw[e] (a) -- node[pos=0.42, above left=0pt, font=\scriptsize, fill=white, inner sep=1.5pt] {1} (b);
        \draw[e] (a) -- node[pos=0.42, above right=0pt, font=\scriptsize, fill=white, inner sep=1.5pt] {1} (c);
        \draw[e] (b) -- node[pos=0.58, left=0pt, font=\scriptsize, fill=white, inner sep=1.5pt] {9} (d);
        \draw[e] (c) -- node[pos=0.58, right=0pt, font=\scriptsize, fill=white, inner sep=1.5pt] {2} (d);
      \end{tikzpicture}
    \end{column}
    \begin{column}{0.64\textwidth}
      \footnotesize
      Capa 1. Desde 1 llegan a 2 y a 3 el par \((0,1)\).\\
      Costo \(1\). Los dos quedan con id \(0\).

      \medskip
      Capa 2. Hacia 4:
      \begin{center}
      \renewcommand{\arraystretch}{1.12}
      \begin{tabular}{@{}l@{\quad}c@{\quad}c@{}}
        desde & par & costo \\
        2, dígito 9 & \((0,9)\) & \(19\) \\
        3, dígito 2 & \((0,2)\) & \(12\)
      \end{tabular}
      \end{center}
      Gana \((0,2)\). El nodo 4 se queda con \(12\).
    \end{column}
  \end{columns}
\end{frame}
